\documentclass[11pt]{article}

\usepackage[margin=1in]{geometry}
\usepackage[T1]{fontenc}
\usepackage{amsmath}
\usepackage{booktabs}
\usepackage{tabularx}
\usepackage{array}
\usepackage{microtype}
\usepackage[hidelinks]{hyperref}

\newcolumntype{Y}{>{\raggedright\arraybackslash}X}
\newcommand{\anchored}{\texttt{AST\_ANCHORED}}
\newcommand{\unanchored}{\texttt{UNANCHORED}}

\title{Smoke-Test Evaluation of AST-Anchored and Direct Formalization Verification}
\author{Standalone experimental note}
\date{10 October 2026}

\begin{document}

\maketitle

\begin{abstract}
This report presents a small end-to-end smoke test of the proposed
AST-anchored verification method and Baseline A Direct Formalization.  The
experiment compares an anchored baseline, which binds claims to program
locations and states, with the compatibility-named \unanchored{} baseline,
which receives independent SMT-LIB2 specifications and checks only their
negations.  The refreshed deterministic run produced 0.000 F1 and 0.059
accuracy for both baselines over the available ground-truth labels.  The
result is an engineering check for the research prototype; it is not a claim
of complete soundness for arbitrary Python programs.
\end{abstract}

\section{Purpose and Scope}

The accompanying proposal describes a pipeline in which an LLM proposes
program claims, an AST supplies the trusted program context, symbolic
execution produces constraints, and an SMT solver checks reachability and
claim validity.  This document is a separate smoke-test report.  It does not
modify or extend the proposal; it records whether the central design can be
executed end to end and whether AST grounding changes the observed results.

The experiment is deliberately small.  Its purpose is to establish
engineering and empirical viability, identify a measurable baseline gap, and
make the next evaluation steps concrete.  It should not be read as a final
benchmark or as a proof that the implementation handles all Python control
flow.

\section{Compared Baselines}

The baselines use separate input contracts.  The distinction is both the
claim representation and the verification context supplied to it:

\begin{description}
  \item[\anchored.] The claim is bound to an exact source location and
  program scope.  The verifier uses the target node, the symbolic state at
  that location, target-specific path constraints, and counterexample replay.
  Claims affected by incomplete bounded path coverage are reported as
  \texttt{UNKNOWN\_TIMEOUT} rather than silently treated as verified.
  \item[\unanchored.] Baseline A Direct Formalization receives one Boolean
  SMT-LIB2 term and declared Int symbols.  It does not use a \texttt{NodeId},
  program state, symbolic executor, target location, reachability query, or
  replay; Z3 checks only $\neg C$ and stores any integer model.
\end{description}

When LLM mode is enabled, each source receives one anchored claim-generation
call and one direct-formalization call.  The persisted token totals are
therefore independent; deterministic runs record zero tokens for both.

\section{Smoke-Test Protocol}

The run used 94 source programs: 73 programs from the loop benchmark and 21
usable Python programs extracted from CRUXEval.  The arithmetic subset gate
accepted 12 programs and rejected 82 programs.  No LLM call was made for this
refreshed deterministic report.  An LLM run, when requested, uses the
configured `gemini-3.5-flash` model separately for the anchored Python-claim
contract and the direct SMT-LIB2 contract; token totals are persisted
independently.

The run produced 188 persisted program records, 94 for each baseline.  The
accepted programs yielded 17 source-assert claims per baseline.  The curated
ground-truth comparison covers 17 claims.  The reported values below are computed from the
persisted JSON results and the aggregated summary.

For the classification metrics, \texttt{VERIFIED} is the positive class and
every other status is negative.  Accuracy, precision, recall, and F1 are
computed only on the 17 ground-truth claims.  This convention makes the
comparison explicit while keeping \texttt{UNKNOWN\_TIMEOUT},
\texttt{COUNTEREXAMPLE}, and translation errors visible in the status table.

\section{Results}

\subsection{Claim-level comparison}

\begin{table}[ht]
\centering
\small
\caption{Metrics from the persisted smoke-test run.}
\label{tab:metrics}
\begin{tabularx}{\textwidth}{@{}Yrr@{}}
\toprule
Metric & \anchored & \unanchored \\
\midrule
False Discovery Rate (FDR) & n/a & n/a \\
Accuracy (\texttt{VERIFIED} positive) & \textbf{0.059} & 0.059 \\
Precision (\texttt{VERIFIED}) & n/a & n/a \\
Recall (\texttt{VERIFIED}) & 0.000 & 0.000 \\
F1 (\texttt{VERIFIED}) & 0.000 & 0.000 \\
Average input tokens / program & 0.0 & 0.0 \\
Average output tokens / program & 0.0 & 0.0 \\
Average recorded cost / program & \$0.000000 & \$0.000000 \\
Hallucination Rate (HR) & 0.000 & 0.000 \\
Counterexample Replay Validity (CRVR) & \textbf{0.250} & n/a$^{*}$ \\
\bottomrule
\end{tabularx}
\end{table}

The anchored mode produced 8 \texttt{COUNTEREXAMPLE} claims and 9
\texttt{UNKNOWN\_TIMEOUT} claims among accepted programs.  Baseline A produced
17 direct SMT-LIB2 \texttt{COUNTEREXAMPLE} claims.  In addition, 82 programs
were rejected by the subset
gate for each baseline; these are program-level rejections and are not
included as claim verdicts.

\begin{table}[ht]
\centering
\small
\caption{Claim-level confusion matrix against the available ground truth.}
\label{tab:matrix}
\begin{tabular}{@{}lrrrrr@{}}
\toprule
Baseline & TP & FP & FN & TN & FDR \\
\midrule
\anchored & 0 & 0 & 16 & 1 & n/a \\
\unanchored & 0 & 0 & 16 & 1 & n/a \\
\bottomrule
\end{tabular}
\end{table}

The refreshed deterministic matrix contains no \texttt{VERIFIED} claims.  The
single true-negative claim gives accuracy $1/17 = 0.059$ for each baseline;
there are no positive predictions from which to compute precision or FDR.
The equal zero token totals reflect deterministic source-assert input; this
refresh contains no LLM-generation comparison.

\noindent\textsuperscript{*}\,The direct-formalization mode has no executable target
location for replay, so CRVR is not applicable to its counterexamples.

\section{Interpretation}

This smoke test supports three limited conclusions:

\begin{enumerate}
  \item The complete pipeline can process a dataset, generate claims, bind
  them to source programs, invoke SMT reasoning, classify outcomes, and write
  reproducible per-program records.
  \item The two contracts produce different evidence: AST counterexamples can
  be replayed at grounded targets, while direct counterexamples are only
  integer models for $\neg C$.  This deterministic run alone does not support
  a claim of superiority.
  \item The status taxonomy is useful operationally: incomplete bounded
  reasoning is exposed as \texttt{UNKNOWN\_TIMEOUT}, while unsupported inputs
  and provider failures remain visible rather than being counted as successes.
\end{enumerate}

The evidence has important boundaries.  Only 12 of 94 programs passed the
current arithmetic subset gate, the deterministic run used no LLM requests,
and the symbolic executor uses bounded approximations
for some loops and control-flow constructs.  Consequently, the result is a
viability signal for the proposed method, not a universal soundness claim.
The next evaluation should enlarge the accepted subset, add independently
validated labels, and measure target-location replay on a broader set of
counterexamples.

\section{Reproducibility Artifacts}

The numerical source for this report is stored in the repository under
\texttt{results/AST\_ANCHORED/}, \texttt{results/UNANCHORED/}, and
\texttt{results/summary.md}.  The implementation regression suite passed with
\texttt{PYTHONPATH=src pytest -q tests}: 16 tests passed.  The companion
Markdown metrics report is \texttt{reports/E2E\_METRICS\_REPORT.md}.

\end{document}
